\section{Proofs of Asymptotic Properties for C-Learner from \Cref{sec:theory_text}}
\label{sec:theory_appendix}
Here, we prove results in \cref{sec:theory_text}. 
We show that C-Learner is semiparametrically efficient (\cref{sec:asymptotics}) and doubly robust (\cref{sec:double_robustness}), under assumptions in \cref{sec:theory_text}. 
We also show that versions of one-step estimation methods (self-normalized AIPW) and targeting methods (TMLE with unbounded continuous outcomes) can satisfy these assumptions (\cref{sec:show_constant_shift}, \cref{sec:show_tmle}), so that semiparametric efficiency and double robustness hold for them immediately as well. 

\subsection{Proof of Theorem~\ref{thm:asymp}}
\label{sec:asymptotics}









Our proof follows a standard argument. We first use the distributional Taylor expansion~\eqref{eq:taylor} to rewrite the estimation error $\what\psi_n^C-\psi$ as the sum of three terms. Then, we address these terms one by one, and we will show how only one of these terms contributes to asymptotic variance. 

Let $Z=(X,A,Y)$ as defined in \cref{sec:ate}. Let $P$ denote the true population distribution of $Z$. 
Let $\psi(P)=P[P[Y\mid X]]$ as in \Cref{sec:ate}.
Functionals $\psi$ may admit a distributional Taylor expansion, also known as a von Mises expansion \citep{fernholz2012mises}, where for any distributions $P,\wb P$ on $Z$, we can write
\begin{align}
\label{eq:taylor}
    \psi(\wb P)-\psi(P)=-\int \varphi(z;\wb P)dP(z)+R_2(\wb P, P)
\end{align}
where 
$\varphi(z;P)$ which can be thought of as a ``gradient'' satisfying the directional derivative formula 
$\frac{\partial}{\partial t}\psi(P+t(\wb P-P))\mid_{t=0}=\int \varphi(z;P)d(\wb P-P)(z)$. (W.l.o.g. we assume $\varphi(z;P)$ is centered so that
$\int \varphi(z;P)dP(z)=0$.)
Here, $-\int \varphi(z;\wb P)dP(z)$ is the first-order term, and $R_2(\wb P, P)$
is the second-order remainder term, which only depends on products or squares of differences between $P,\wb P$.

When $\psi$ is the ATE as in our setting, $\psi$ admits such an expansion (e.g. using Theorem 20.8 in \citet{VanDerVaart98}; $\psi(P)$ is Hadamard differentiable as it is linear in $P$. Also see \citep{Kennedy22} for a primer.).
In such an expansion, $\varphi$ is as calculated in \citep{Hahn98} (also see \citep{Kennedy22}) as
\begin{align}
\label{eq:full_if}
\varphi(Z;\wb P):=\frac{A}{\wb \pi(X)}(Y-\wb \mu(X))+\wb \mu(X)-\psi(\wb P),
\end{align}
where $\wb\pi(x):=\wb P[A=1\mid X]$
and $\wb \mu(x):=\wb P[Y\mid A=1,X]$. In particular, we get the following explicit formula for the second-order term
\begin{align}
R_2(\wb P,P) := \int \pi(x) \left( \frac{1}{\wb\pi(x)} - \frac{1}{\pi(x)} \right) \left( \wb \mu (x) - \mu(x) \right) dP(x).
\end{align}
We will apply this to our C-Learner estimator as defined in \cref{sec:theory_text}. Recall that
$\what \pi_{-k,n}$ is trained to predict treatment $A$ given $X$ using $P_{-k,n}$, and $\what \mu_{-k,n}^C$ is trained to predict outcome given $X$ and $A=1$ on $P_{-k,n}$, under the constraint that $P_{k,n}\left[\frac{A}{\what \pi_{-k,n}(X)}(Y-\what \mu_{-k,n}^C(X)\right]=0$.
Our C-Learner estimator is the mean of plug-in estimators across folds: for each fold, write $\what\psi^C_{k,n}=\psi(\what P_{k,n}^C)$ so the C-Learner estimate is the average $\what\psi^C_n=\frac{1}{K}\sum_{k=1}^K \what\psi^C_{k,n}$. 


Noting that any distribution decomposes $\wb{P}=\wb{P}_X\times \wb{P}_{A\mid X}\times \wb{P}_{Y\mid A,X}$ and $\psi(\wb{P})=\wb{P}_X[\mu(X)]$, the following definitions
\begin{align*}
\what P^C_{X;k,n}&:=P_{X;k,n} \\
\what P^C_{A\mid X ;k,n}[A=1\mid X=x]&:=\what \pi_{-k,n}(x) \\ 
\what P^C_{Y\mid A,X ;k,n}[Y\mid A=1, X=x]&:=\what \mu_{-k,n}^C(x)
\end{align*}
provide a well-defined joint distribution  $\what P^C_{k,n}$.

For each data fold $k$, we use the distributional Taylor expansion above, where we replace $\wb P$ with the joint distribution $\what P^C_{k,n}$
\begin{align}
\label{eq:taylor_c_learner}
\psi(\what P^C_{k,n}) - \psi(P) &= - P \varphi(Z;\what P^C_{k,n}) + R_2(\what P^C_{k,n},P) \nonumber\\ 
&= (P_{k,n}-P)\varphi(Z;P)\nonumber
-P_{k,n} \varphi(Z;\what P^C_{k,n})
\\&\quad  +(P_{k,n}-P)
(\varphi(Z;\what P^C_{k,n})-\varphi(Z;P))
+R_2(\what P^C_{k,n},P). 
\end{align}
Observe that by using Equation~\eqref{eq:full_if} and the definition of the C-Learner,
\begin{align*}
    P_{k,n}\varphi(Z;\what P^C_{k,n})
&=P_{k,n}\left[\frac{A}{\what\pi_{-k,n}(X)}(Y-\what \mu^C_{-k,n}(X))+\what\mu^C_{-k,n}(X)-\psi(\what P^C_{k,n})\right]
\\&=\underbrace{P_{k,n}\left[\frac{A}{\what\pi_{-k,n}(X)}(Y-\what \mu^C_{-k,n}(X))\right]}_{=0 \text{ by C-Learner constraint}}
+\underbrace{P_{k,n}[\what\mu^C_{-k,n}]-P_{k,n}[\what\mu^C_{-k,n}]}_{=0}=0.
\end{align*}
Taking the average of Equation~\eqref{eq:taylor_c_learner} over $k=1,\ldots,K$,
we can write the error $\what\psi^C_n-\psi$ as the sum of three terms.
    \begin{align}   
     \what \psi^C_{n} - \psi  
    & = \frac{1}{K}\sum_{k=1}^K\underbrace{(P_{k,n}-P)\varphi(Z;P)}_{S^*_k} \nonumber \\
    &  \qquad + \frac{1}{K}\sum_{k=1}^K\underbrace{(P_{k,n}-P) \left(\varphi(Z;\what P^C_{k,n})-\varphi(Z;P)\right)}_{T_{1k}} + \frac{1}{K}\sum_{k=1}^K\underbrace{R_2(\what P^C_{k,n},P)}_{T_{2k}}. \label{eq:S+T1+T2}
\end{align}
Using the decomposition~\eqref{eq:S+T1+T2}, we write 
\begin{align}
S^*=\frac{1}{K}\sum_{k=1}^K S^*_k,  \hspace{1em}
T_1=\frac{1}{K}\sum_{k=1}^K T_{1k},  \hspace{1em}
T_2=\frac{1}{K}\sum_{k=1}^K T_{2k},
\end{align}
so that
$\what \psi^C_{n} - \psi = S^* + T_1 + T_2$.
We address the terms $S^*,T_1$ and $T_2$ separately. The first term can be rewritten as
$$
S^* = \frac{1}{K}\sum_{i=1}^K (P_{k,n}-P)\varphi(Z;P) = (P_n-P)\varphi(Z;P)
$$
so that by the central limit theorem, 
$$
\sqrt{n}S^* \cd {N}\left(0, \var_P(\varphi(Z;P))\right).
$$ 
Observe that this quantity depends only on $\psi$ and $P$, so that it cannot be made smaller by choice of estimator. 
If the variance of an estimator for $\psi$ is $\var_P(\varphi(Z;P))$, then it is semiparametrically efficient in the local asymptotic minimax sense (Theorem 25.21 of \cite{van2000asymptotic}). 
Thus,  it suffices to show that the rest of the terms, $T_1$ and $T_2$, are $o_P(n^{-1/2})$, so that
$$
\sqrt{n}(\what \psi^C_n - \psi)=\sqrt{n} S^*+o_P(1)\cd {N}\left(0, \var_P (\varphi(Z;P))\right)
.$$
For a fixed $k$,
$|T_{1k}|=o_P(n^{-1/2})$ 
by Assumption~\ref{ass:empirical_proc_mean}
so that $|T_1|=o_P(n^{-1/2})$ as desired. 
The second-order remainder term in the distributional Taylor expansion \eqref{eq:taylor} where
we replace $\wb P$ with $\what{P}^C_{k,n}$
is
$$
T_{2k} = \int \pi(x) \left( \frac{1}{\what\pi_{-k,n}(x)} - \frac{1}{\pi(x)} \right) \left( \what\mu^C_{-k,n}(x) - \mu(x) \right) \, dP(x)
.$$
Under the overlap assumption (Assumption~\ref{ass:overlap}) and  Cauchy-Schwarz, 
\begin{align}
    |T_{2k}| &\leq \frac{1}{\eta} \int \left| \what\pi_{-k,n}(x) - \pi(x) \right| \left| \what\mu^C_{-k,n}(x) - \mu(x) \right| \, dP(x) \nonumber \\
    &\leq \frac{1}{\eta} \left\|\what \pi_{-k,n} -\pi\right\|_{L_2(P)} \left\|\what \mu^C_{-k,n} - \mu\right\|_{L_2(P)}. \label{eq:T2_cauchy_schwarz}
\end{align}
By Assumption~\ref{ass:consistency},
$|T_{2k}|=o_P(n^{-1/2})$
so that 
$|T_{2}|=o_P(n^{-1/2})$ as desired. 





\subsection{Proof of Theorem~\ref{thm:dr}}
\label{sec:double_robustness}
We briefly sketch the proof as it is a minor modification of our previous proof in \cref{sec:asymptotics}.  To show the C-Learner estimator 
$\what\psi^C_n$ 
is consistent (rather than that our estimator has the desired asymptotics) under Assumption~\ref{ass:risk_decay} (rather than Assumption~\ref{ass:consistency}), we  again rewrite the error $\what\psi_n^C-\psi$ as a sum of three terms and show  each converges to 0 in probability: 
\begin{itemize}
    \item $S^*$: This term converges to 0 in probability. 
    \item $T_{2k}$: This also converges to 0 in probability by the same logic as before. Note we only need this to be $o_P(1)$ and not $o_P(n^{-1/2})$ as we would have required for efficiency. 
    \item $T_{1k}$: This converges to 0 in probability by Assumption~\ref{ass:empirical_proc_mean}. 
\end{itemize}

\subsection{Showing Self-Normalized AIPW and TMLE Satisfy C-Learner Conditions for Theorems~\ref{thm:asymp} and \ref{thm:dr}}
\label{sec:show_for_aipw_tmle}

As self-normalized AIPW and TMLE involve simple adjustments to the unconstrained outcome models, 
in this section, we state assumptions on the \emph{unconstrained} outcome model
$\what{\mu}_{-k, n}$, fitted in the usual manner on the auxiliary fold $P_{-k, n}$
\begin{equation*}
  \what{\mu}_{-k, n} \in \argmin_{\wt{\mu} \in \mc{F}} 
  P_{-k, n}[ A (Y - \wt{\mu}(X))^2],
\end{equation*}
and also assumptions on the gap between the constrained and unconstrained outcome models. Also in this section, we show how these aforementioned assumptions satisfy the C-Learner assumptions required for Theorems~\ref{thm:asymp} and \ref{thm:dr}. 

Then in the following sections, we show how
self-normalized AIPW and TMLE satisfy the assumptions stated in this section on the gap between the constrained and unconstrained outcome models. 
 
\paragraph{Unconstrained Outcome Models}
The following assumption on unconstrained outcome models is analogous to Assumption~\ref{ass:consistency}. This assumption is standard. 
\begin{assumption}[Convergence rates of propensity and \emph{unconstrained} outcome models]
\label{ass:consistency_unconstrained}
For all~$k=1,\ldots,K$,
    $$
    \|\what\pi_{-k,n} - \pi\|_{L_2(P)}=o_P(n^{-\frac{1}{4}}), \quad  \|\what\mu_{-k,n}^C - \mu\|_{L_2(P)}=o_P(n^{-\frac{1}{4}}).
    $$
\end{assumption}

\paragraph{Distance Between Constrained and Unconstrained Outcome Models}
These assumptions essentially ensure that the first-order constraint~\eqref{eq:constraint} does not change the outcome model too much asymptotically, i.e., $\what{\mu}^C_{-k,n}$ and $\what{\mu}_{-k,n}$ are similar asymptotically. 
This first constraint is used to show
Assumption~\ref{ass:consistency}:
\begin{assumption}[The constraint is negligible asymptotically] For all $k$ = 1, \ldots, K,
\label{ass:muC_vs_mu}
$$\|\what\mu^C_{-k,n}-\what\mu_{-k,n}\|_{L_2(P)}=o_P(n^{-1/4}).$$
\end{assumption}
Notably, the distance between the constrained solution $\what\mu^C_{-k,n}$ and its unconstrained counterpart $\what\mu_{-k,n}$ can be as large as that between $\what\mu_{-k,n}$ and the true parameter $\mu$, asymptotically.

This next assumption is used to show Assumption~\ref{ass:empirical_proc_mean}:
\begin{assumption}[Empirical process assumption on $\what\mu^C_{-k,n}$ vs $\what\mu_{-k,n}$]
\label{ass:empirical_proc_mean_diff}\
$$(P_{k,n}-P)(  \what\mu_{-k,n}^C(X)- \what\mu_{-k,n}(X)) =o_P(n^{-1/2}).$$
\end{assumption}
Given these assumptions, we show the C-Learner assumptions hold. 
\begin{proposition}[C-Learner conditions hold, given assumptions on unconstrained outcome models and the gap between constrained and unconstrained models]\
\label{prop:translate}

Assume Assumptions~\ref{ass:overlap}, 
\ref{ass:bounded_outcomes},
\ref{ass:consistency_unconstrained},
and \ref{ass:empirical_proc_mean_diff}.
Then Assumptions~\ref{ass:overlap}, 
\ref{ass:bounded_outcomes},
\ref{ass:consistency}, \ref{ass:empirical_proc_mean}
are satisfied. 
\end{proposition}
To show this proposition, it suffices to show Assumptions~\ref{ass:consistency} and \ref{ass:empirical_proc_mean}. We will show these assumptions in the rest of this section.
\paragraph{Showing  Assumption~\ref{ass:consistency}}
By the triangle inequality, 
\begin{align}
    \|\what \mu^C_{-k,n}-\mu\|_{L_2(P)} &\leq \|\what \mu_{-k,n}-\mu\|_{L_2(P)} + \|\what \mu^C_{-k,n}-\what\mu_{-k,n} \|_{L_2(P)}
.\end{align}

\paragraph{Showing  Assumption~\ref{ass:empirical_proc_mean}}
Applying the triangle inequality again yields
\begin{align}
|T_{1k}|&:=\left|(P_{k,n}-P)\left(\varphi(Z;\what P_{k,n}^C) -\varphi(Z; P)\right)\right| \nonumber
\\&\leq \left| (P_{k,n}-P)\left( \varphi(Z;\what P_{k,n}^C) -\varphi(Z; \what P_{-k,n}) \right) \right| + 
 \left| (P_{k,n}-P)\left( \varphi(Z; \what P_{-k,n}) -\varphi(Z; P) \right) \right|
 \label{eq:T1k_bound}
\end{align}
where $\what P_{k,n}$ is defined as
\begin{align*}
\what P_{X ; k,n}&:=P_{X;k,n}\\ 
\what P_{A\mid X ; k,n}[A=1\mid X=x]&:=\what \pi_{-k,n}(x) \\ 
\what P_{Y\mid A,X ; k,n}[Y\mid A=1, X=x]&:=\what \mu_{-k,n}(x)
\end{align*}
with $\what \pi_{-k,n},\what \mu_{-k,n}$ 
as defined in \cref{sec:theory_text},
and $\what P_{k,n}^C$ is defined as in \cref{sec:theory_text}. 
The second term in~\eqref{eq:T1k_bound} is addressed 
using a standard argument for cross-fitting, as $P_{k,n}$ and $\what P_{-k,n}$ (and therefore $\varphi(Z;\what P_{-k,n})$) use disjoint data. In contrast, the first term is not handled by standard cross-fitting arguments: $\what P_{-k,n}$ only uses data from all but the $k$-th fold, while $\what P^C_{k,n}$ uses the $k$-th fold, except that $\what \mu^C_{-k,n}$ is made to satisfy a constraint that \emph{does} use the $k$-th fold, as described in \cref{sec:theory_text}. We begin by addressing the second term, which is more standard. %

\begin{lemma}[Cross-fitting lemma]
\label{lem:cross_fitting}
Let $\what{f}(z)$ be a function estimated from an iid sample $Z^N=$ $\left(Z_{n+1}, \ldots, Z_N\right)$, and let ${P}_n$ denote the empirical measure over $\left(Z_1, \ldots, Z_n\right)$, which is independent of $Z^N$. Let $f$ be the function estimated from the full distribution $P$. Then (omitting arguments for brevity)
$
\left({P}_n-{P}\right)(\what{f}-f)$
has zero $P$-expectation, and $P$-variance upper bounded by
$\frac{1}{n}\|\what{f}-f\|_{L_2(P)}^2.$
\end{lemma}
\begin{proof}
First note that the conditional mean is 0, i.e. $P\left[(P_n-P)(\what f-f )\mid Z^N\right]=0$ since
$$
P\left[{P}_n(\what{f}-f) \mid Z^N\right]=P\left(\what{f}-f \mid Z^N\right)={P}(\what{f}-f) .
$$
The conditional variance is
\begin{align*}
    \var_P\left\{\left({P}_n-{P}\right)(\what{f}-f) \mid Z^N\right\} &=\var_P\left\{{P}_n(\what{f}-f) \mid Z^N\right\}\\
    &=\frac{1}{n} \var_P\left(\what{f}-f \mid Z^N\right) \\
    &\leq \frac{1}{n}\|\what{f}-f\|_{L_2(P)}^2.
\end{align*}
Then for (unconditional) mean and variance, $P\left[(P_n-P)(\what f-f)\right]=0$ and
\begin{align*}
&\var_P\left\{(P_n-P)(\what f-f)\right\} \\
&= 
\var_P \left\{P\left[(P_n-P)(\what f-f)\mid Z^N\right]  \right\}  
+P\left[\var_P\left\{(P_n-P)(\what f-f)\mid Z^N\right\}\right]
\\
&\leq 0+\frac{1}{n}\|\what{f}-f\|_{L_2(P)}^2.
\end{align*}
\end{proof}
Now we show the second term in the RHS in Equation~\eqref{eq:T1k_bound} is $o_P(n^{-1/2})$. To do this, write
\begin{align}
\varphi(Z;P)&=\frac{A}{\pi(X)}(Y-\mu(X))+\mu(X)-\psi(P)\\
\varphi(Z;\what{P}_{-k,n})&=\frac{A}{\what\pi_{-k,n}(X)}(Y-\what\mu_{-k,n}(X))+\what\mu_{-k,n}(X)-\psi(\what{P}_{-k,n})    
.\end{align}
Observe that $(P_{k,n}-P)\psi(P)=(P_{k,n}-P)\psi(\what {P}_{-k,n})=0$ as $\psi(P),\psi(\what P_{-k,n})$ are constants.
Thus, 
it remains to show that for a fixed $k$,
$(P_{k,n}-P)( \what f_{k,n} - f)=o_P(n^{-1/2})$, 
where we omit arguments for brevity and let 
\begin{align}
\label{eq:def_f_in_emp_proc}
    \what f_{k,n} &= \frac{A}{\what \pi_{-k,n}}(Y-\what \mu_{-k,n})+\mu_{-k,n}, \\
    f &= \frac{A}{\pi}(Y-\mu) +\mu.
\end{align}
We do this by using Lemma~\ref{lem:cross_fitting}. 
Observe that
\begin{align}
\label{eq:second_term_t1}
    \what f_{k,n} - f =\left(1+\frac{A}{\pi}\right)(\mu-\what\mu_{-k,n})
    +\frac{A}{\what\pi_{-k,n}\cdot  \pi} (Y-\what\mu_{-k,n})(\pi-\what\pi_{-k,n}).
\end{align}
Then using Assumption~\ref{ass:overlap},
\begin{align}
\label{eq:T1_f_both}
    \|\what f_{k,n} - f\|_{L_2(P)}
    &\leq 
    \left(1+\frac{1}{\eta}\right) \|\what\mu_{-k,n}- \mu\|_{L_2(P)}
    +\frac{1}{\eta^2} \| A(Y-\what\mu_{-k,n})\|_{L_2(P)} \|\pi-\what\pi_{-k,n}\|_{L_2(P)}.
\end{align}
The leftmost term on the RHS converges to 0,  
using Assumption~\ref{ass:consistency_unconstrained}. 
Note that we can bound the rightmost term using the triangle inequality
\begin{align}
\label{eq:T1_f_second}
    \left(1+\frac{1}{\eta^2}\right) \left(\| A(Y-\mu)\|_{L_2(P)} +  \| \mu - \what\mu_{-k,n}\|_{L_2(P)} \right) \|\pi-\what\pi_{-k,n}\|_{L_2(P)},
\end{align}
which also converges to 0 by %
Assumption~\ref{ass:consistency_unconstrained}
and Assumption~\ref{ass:bounded_outcomes}. 
Thus combining with Lemma~\ref{lem:cross_fitting} we obtain that the second term on the RHS of Equation~\eqref{eq:T1k_bound} is $o_P(n^{-1/2})$,
as $n^{1/2}(P_{k,n}-P)( \what f_{k,n} - f)$ has mean 0 and variance $\leq \|\what f_{k,n} - f\|_{L_2(P)}$ which converges to 0 as $n\to\infty$.

Now we address the remaining term: we show the first term on the RHS of Equation~\eqref{eq:T1k_bound} is $o_P(n^{-1/2})$. We use a similar argument to before: 
let $\what f_{k,n}$ be as before, and $\what f^C_{k,n}$ as below:
\begin{align}
    \what f_{k,n} &:= \frac{A}{\what \pi_{-k,n}}(Y-\what \mu_{-k,n})+\what \mu_{-k,n}\\
    \what f^C_{k,n} &:= \frac{A}{\what \pi_{-k,n}}(Y-\what \mu_{-k,n}^C)+\what \mu_{-k,n}^C
\end{align}
where we again omit the $\psi(\what P^C_{-k,n}),\psi(\what P_{-k,n})$ terms in 
$\varphi(\what P_{k,n}^C),\varphi(\what P_{k,n})$
from $\what f^C_{k,n},\what f_{k,n}$ 
above as they are constants. 
We can't use Lemma~\ref{lem:cross_fitting}
since $\what f^C_{k,n}$ also uses $P_{k,n}$ to fit, so instead, by 
using Assumptions \ref{ass:overlap} and then \ref{ass:empirical_proc_mean_diff},
\begin{align}
    (P_{k,n}-P)
\left(\varphi(Z;\what P^C_{k,n})-\varphi(Z;\what P_{k,n})\right)
&=
(P_{k,n}-P)\left(\what f^C_{k,n}-\what f_{k,n}\right)
\\&\leq \left(1+\frac{1}{\eta}\right)(P_{k,n}-P)(\what\mu^C_{-k,n}(X)-\what\mu_{-k,n}(X))\\&=o_P(n^{-1/2}).
\end{align}


\subsection{Self-Normalized AIPW Satisfies Assumptions~\ref{ass:muC_vs_mu} and \ref{ass:empirical_proc_mean_diff}}
\label{sec:show_constant_shift}
Recall that we showed how self-normalized AIPW also satisfies the C-Learner formulation in \cref{sec:other_methods}. 
Here we show that Assumptions~\ref{ass:muC_vs_mu} and \ref{ass:empirical_proc_mean_diff} hold for self-normalized 
AIPW, so that Theorems~\ref{thm:asymp} and \ref{thm:dr} follow through Proposition~\ref{prop:translate}. 



As in the discussion in \cref{sec:other_methods}, self-normalized AIPW is equivalent to the specific C-Learner $\what\mu^C_{-k,n}$ that is defined by adjusting $\what\mu_{-k,n}$ by an additive constant:
$$\what\mu^C_{-k,n}(x)=\what\mu_{-k,n}(x)+c_{k,n}\;\text{ where }\;c_{k,n}:=\frac{P_{k,n}\left[\frac{A}{\what\pi_{-k,n}(X)}(Y-\what\mu_{-k,n}(X))\right]}{P_{k,n}\left[\frac{A}{\what\pi_{-k,n}(X)}\right]}.$$

\paragraph{Showing Assumption~\ref{ass:muC_vs_mu}:}
Since $\what\mu_{-k,n}^C$
is just a constant offset from $\what\mu_{-k,n}$, it suffices to show 
that $c_{k,n}=o_P(n^{-1/4})$, for a fixed $k$. 

First, we address the denominator of $c_{k,n}$. Let
$c_{k,n}^{\rm den}:=P_{k,n}\left[\frac{A}{\what\pi_{-k,n}(X)}\right]$ 
and we will show $1/c_{k,n}^{\rm den}=O_P(1)$. 
First note $c_{k,n}^{\rm den}=P_{k,n}\left[\frac{A}{\what\pi_{-k,n}(X)}\right] \geq P_{k,n}[A]$. %
Then observe that $P_{k,n}[A]\overset{p}{\to} P[A]$, and that $1/P_{k,n}[A]<\infty$ a.s. for large enough $n$ by Borel-Cantelli lemma (as the probability of the event that $P_{k,n}[A]=0$ is finitely-summable, as $P(A=0)<1$ by the overlap assumption (Assumption~\ref{ass:overlap}). %
Then $1/P_{k,n}[A]\overset{p}{\to} 1/P[A]$ for large enough $n$, 
so that $1/P_{k,n}[A]=O_P(1)$. %

Now we address the numerator of $c_{k,n}$. 
For brevity, call this $c_{k,n}^{\text{num}}$.
\begin{align}
c_{k,n}^{\text{num}} &=    P_{k,n}\left[\frac{A}{\what\pi_{-k,n}(X)}(Y-\what\mu_{-k,n}(X))\right] 
\\&= P\left[\frac{A}{\what\pi_{-k,n}(X)}(Y-\what\mu_{-k,n}(X))\right] 
+(P_{k,n}-P)\left[\frac{A}{\what\pi_{-k,n}(X)}(Y-\what\mu_{-k,n}(X))\right]
\\&= \underbrace{P\left[\frac{A}{\what\pi_{-k,n}(X)}(Y-\mu(X))\right]}_{=0}
+P\left[\frac{A}{\what\pi_{-k,n}(X)}(\mu(X)-\what\mu_{-k,n}(X))\right]\\&\quad 
+(P_{k,n}-P)\left[\frac{A}{\what\pi_{-k,n}(X)}(Y-\what\mu_{-k,n}(X))\right]\nonumber
\\&=P\left[\frac{A}{\what\pi_{-k,n}(X)}(\mu(X)-\what\mu_{-k,n}(X))\right]
+(P_{k,n}-P)\left[\frac{A}{\what\pi_{-k,n}(X)}(Y-\what\mu_{-k,n}(X))\right]
\end{align}
so that
\begin{align}
    |c_{k,n}^{\rm num}| & \leq %
    \frac{1}{\eta} P\left|\mu(X)-\what\mu_{-k,n}(X)\right|
    +\frac{1}{\eta} (P_{k,n}-P)\left|Y-\what\mu_{-k,n}(X)\right|
\end{align}
where the inequality is by the overlap assumption (Assumption~\ref{ass:overlap}) as $A/\what\pi_{-k,n}(X)\leq 1/\eta$. 
We show the terms in the last line are all $o_P(n^{-1/4})$. 
The first term is upper bounded by $\frac{1}{\eta} \|\what\mu_{-k,n}(X)-\mu(X)\|_{L_2(P)}=o_P(n^{-1/4})$ by Assumption~\ref{ass:consistency_unconstrained}. 
For the second term, 
we use Lemma~\ref{lem:cross_fitting} to show
$n^{1/4}(P_{k,n}-P)|Y-\what\mu_{-k,n}(X)|\overset{p}{\to}0$:
it has mean of 0 and variance
$\leq \frac{n^{1/2}}{n}\|Y-\what\mu_{-k,n}\|_{L_2(P)}$. 
This upper bound on variance goes to 0, %
which follows from Assumptions~\ref{ass:consistency_unconstrained} and \ref{ass:bounded_outcomes}:
\begin{align*}
    \|Y-\what\mu_{-k,n}\|_{L_2(P)} &\leq \|Y-\mu\|_{L_2(P)} + \|A(\mu-\what\mu_{-k,n})\|_{L_2(P)} \\
    &\leq B + \|\mu-\what\mu_{-k,n}\|_{L_2(P)} \\
    &= B + o_P(n^{-1/4}).
\end{align*}
We arrive at the desired result $c_{k,n}=o_P(n^{-1/4})$ as 
$c_{k,n}=c_{k,n}^{\rm num}/c_{k,n}^{\rm den}$, and we just showed that 
$1/c_{k,n}^{\rm den}=O_P(1)$ and 
$c_{k,n}^{\rm num}=o_P(n^{-1/4})$. 
Note that we have shown Assumption~\ref{ass:muC_vs_mu} for any sequence of $\what\pi_{-k,n}$'s, as they are bounded by a constant. 


\paragraph{Showing Assumption~\ref{ass:empirical_proc_mean_diff}:}
To show $(P_{k,n}-P)(  \what\mu_{-k,n}^C(X)- \what\mu_{-k,n}(X)) =o_P(n^{-1/2})$, 
in the case of the constant shift $\what\mu^C(X)=\what\mu(X)+c_{k,n}$, 
note that  $\what \mu^C_{-k,n}(X)-\what \mu_{-k,n}(X) = c_{k,n}$ for every $X$. Therefore, $$
(P_{k,n}-P)(  \what\mu_{-k,n}^C(X)- \what\mu_{-k,n}(X)) = (P_{k,n}-P)c_{k,n} = 0,
$$
with $c_{k,n} < \infty$ a.s. for large enough $n$ (which exists by Borel-Cantelli lemma, as in when we showed Assumption~\ref{ass:muC_vs_mu}). 



\subsection{TMLE Satisfies Assumptions~\ref{ass:muC_vs_mu} and \ref{ass:empirical_proc_mean_diff}}
\label{sec:show_tmle}
Recall that we showed how a version of the TMLE for estimating the ATE with continuous unbounded outcomes also satisfies the C-Learner formulation in \cref{sec:other_methods}. 
Here we show that a cross-fitted version of the TMLE for estimating the ATE with continuous unbounded outcomes additionally satisfies Assumptions~\ref{ass:muC_vs_mu} and \ref{ass:empirical_proc_mean_diff}, so that Theorems~\ref{thm:asymp} and \ref{thm:dr} follow through Proposition~\ref{prop:translate}. 
Note that the formulation below fits a separate $\epsilon^\star_{k,n}$ per cross-fitting split for consistency with C-Learner, rather than one $\epsilon^\star$ overall as described in the cross-validated version of TMLE in \citep{VanDerLaanRo11}. 
$$
\what \mu^{C}_{-k,n}(X) = \what \mu_{-k,n}(X) + \epsilon^\star_{k,n}\frac{A}{\what \pi_{-k,n}(X)}, \; \text{ where } \epsilon^\star_{k,n} = \frac{P_{k,n}\left[\frac{A}{\what\pi_{-k,n}(X)}(Y-\what\mu_{-k,n}(X))\right]}{P_{k,n}\left[\frac{A}{\what\pi_{-k,n}^2(X)}\right]}.
$$
\paragraph{Showing Assumption~\ref{ass:muC_vs_mu}:}
Observe that by Assumption~\ref{ass:overlap}
$$
\left\|\epsilon_{k,n}^\star\frac{A}{\what\pi_{-k,n}(X)}\right\|_{L_2(P)} \leq \frac{1}{\eta}\|\epsilon_{k,n}^\star\|_{L_2(P)}
$$
so it suffices to show that $\epsilon_{k,n}^\star=o_P(n^{-1/4})$.
Note that since $\what\pi_{-k,n}(x) \leq 1$ for all $x$,
\begin{align}
\label{eq:eps}    
|\epsilon_{k,n}^\star| \leq %
\frac{P_{k,n}\left[\frac{A}{\what\pi_{-k,n}(X)}(Y-\what\mu_{-k,n}(X))\right]}{P_{k,n}\left[\frac{A}{\what\pi_{-k,n}(X)}\right]}
\leq |c_{k,n}|
\end{align}
where $c_{k,n}$ is the constant adjustment in \cref{sec:show_constant_shift}, 
so that $\epsilon^\star_{k,n}=o_P(n^{-1/4})$. 


\paragraph{Showing Assumption~\ref{ass:empirical_proc_mean_diff}:}
We want to show
$$
(P_{k,n}-P)(  \what\mu_{-k,n}^\star(X)- \what\mu_{-k,n}(X)) = 
(P_{k,n}-P) \left(\epsilon^\star_{k,n}\frac{A}{\what \pi_{-k,n}(X)}\right) = o_P(n^{-1/2}).
$$
Note that
$$(P_{k,n}-P) \left(\epsilon^\star_{k,n}\frac{A}{\what \pi_{-k,n}(X)}\right)
=\epsilon^\star_{k,n}(P_{k,n}-P)
\left(\frac{A}{\what \pi_{-k,n}(X)}\right)
$$
and that $(P_{k,n}-P)
\left(\frac{A}{\what \pi_{-k,n}(X)}\right)=O_P(n^{-1/2})$ by Lemma~\ref{lem:cross_fitting},
since $|A/\what\pi_{-k,n}(X)|\leq 1/\eta$ a.s. by Assumption~\ref{ass:overlap}. 
Additionally, $\epsilon^\star_{k,n}=o_P(n^{-1/4})$, by the argument above where we showed  Assumption~\ref{ass:muC_vs_mu}. 
The desired result follows by taking the product of the rates for $\epsilon^\star_{k,n}$ and $(P_{k,n}-P)
\left(\frac{A}{\what \pi_{-k,n}(X)}\right)$. %




